% Modernised reading — fonds Alexandre Grothendieck, shelfmark 160,
% pages 1-22 (the whole folder).
%
% This is an INTERPRETATION, not a transcription. It re-reads the pages in
% today's notation and vocabulary, states the hypotheses the manuscript leaves
% implicit, and drops the critical apparatus so that the mathematics reads
% straight through. Everything the brackets and underlines carried — a doubtful
% reading, a notation he used differently, a missing passage — moves into a
% footnote.
%
% It is held to being mathematically correct as it stands, not merely faithful
% to the page. Where the manuscript is loose or mistaken, the true statement is
% what appears here, and the footnote says what the page has. In any dispute of
% substance the transcription governs: whoever cites Grothendieck must cite
% batch-01.fr.tex and batch-02.fr.tex.
%
% Pass: Opus 5 (claude-opus-5), 2026-09-19 — first pass, unchecked against the
% pages by a human. The model was named explicitly in the request; it is the
% model this reading actually ran on.
%
% Scope: a SINGLE document for the whole folder. The shelfmark holds two
% batches — pages 1-20 and pages 21-22 — read together in one pass, and no
% part of the folder is read in another file. Page 17 is blank and carries no
% mathematics; it is absent from the transcription and from here.
%
% Notation fixed for the whole document, and departing from the pages:
%   - A^n for affine n-space, where the pages write E^n; G_m x A^{n-1} for the
%     open subscheme the pages write E^* x E^{n-1}.
%   - X^{(m)} for the m-th infinitesimal neighbourhood of Z in X on page 1,
%     where the page writes X_m. From page 2 on, X_n means something else
%     entirely — the scheme of subgroups of rank p^n — and the collision is
%     the one real notational trap in the folder.
%   - alpha_{p^r} = Ker(F^r), following the pages themselves from page 2b on;
%     NOT alpha_p^r, which is a different group scheme. See the footnote at
%     the surjectivity argument.
%   - Delta_n for the Moore determinant, A_i for the Dickson invariants,
%     indexed as on pages 9-16 (A_i homogeneous of degree p^n - p^i). Page 19
%     indexes them the other way round; that is recorded, not adopted.
%   - V(E) = Spec Sym(E), so that sections of V(E) are sections of the DUAL
%     E-check. This is EGA II 1.7.8 and it is also what pages 13-14 use; the
%     variance matters and is stated in the conventions.
%
% Substantive departures from the pages, each footnoted where it occurs:
%   - page 2: the composite is Frob^r o f_alpha, not f_alpha o Frob^r.
%   - page 4: the special fibre's kernel has rank p^m with m the LAST nonzero
%     index, not p^{n-i} with i the first.
%   - page 9: Delta_0 = (-1)^n Delta_n^p exactly, sign included; this is
%     consistent with A_0 = (-1)^n Delta_n^{p-1} and the two are not in
%     conflict.
%   - page 11: A_i(X_1,...,X_{n-1},0) = A_{i-1}(X_1,...,X_{n-1})^p — the page
%     omits the Frobenius twist.
%   - page 21, right column: the first relation drops the term b(alpha-beta),
%     and the three lines after it are the identity for b*alpha, not b.
%
% What the folder announces and never establishes: the reducedness of X_n in
% general (page 5 asks, pages 5-6 answer only at n = 1 and only at the
% origin); the constant c in the Moore factorisation (page 7 says "reste a
% determiner c"); the descent of page 18; the question mark on the right-hand
% column of page 16 and on the last diagram of page 22.
\documentclass[11pt,a4paper]{article}
% Le préambule commun à toutes les transcriptions du fonds Grothendieck.
%
% Il tient en une idée : **l'incertitude se compose, elle ne se résout pas.**
% Une transcription qui lisse un mot douteux a détruit la seule chose pour
% laquelle on la faisait — un lecteur doit pouvoir savoir, à chaque endroit,
% ce qui était sur le feuillet et ce qui a été ajouté. Les six macros qui
% suivent sont donc l'appareil critique, et non de la décoration.
%
% Les mêmes commandes sont reconnues par `scripts/render.mjs`, qui produit la
% vue de lecture du site. Une macro ajoutée ici doit l'être aussi là-bas,
% faute de quoi le rendu échouera bruyamment — ce qui est voulu : un rendu
% silencieusement faux est pire qu'un rendu absent.

\ProvidesPackage{grothendieck}[2026/08/08 Transcriptions du fonds Grothendieck]

\usepackage{amsmath,amssymb,amsthm}
\usepackage{mathrsfs}
\usepackage{stmaryrd}
% Les diagrammes commutatifs sont partout dans ce fonds : c'est la langue
% même de la géométrie algébrique de Grothendieck, et une transcription qui
% les rendrait en prose ne transcrirait rien.
\usepackage{tikz-cd}
% Français par défaut — c'est la langue des feuillets — mais l'anglais est
% chargé aussi : la traduction et le résumé sont des documents anglais, et la
% typographie française (espaces avant les deux-points) y serait une faute.
% Ces documents basculent par \selectlanguage{english} en tête de corps.
\usepackage[english,french]{babel}
\usepackage{fontspec}
\usepackage{geometry}
\geometry{a4paper,margin=25mm,marginparwidth=28mm}
\usepackage{marginnote}
\usepackage{xcolor}
% ulem, not soul. `soul`'s \st and \ul are built on a character-by-character
% scan that XeLaTeX's font machinery breaks: the document fails with an
% undefined control sequence rather than degrading. ulem does the same two jobs
% and is Unicode-engine safe. `normalem` stops it hijacking \emph, which the
% transcriptions use for Grothendieck's own emphasis.
\usepackage[normalem]{ulem}
\usepackage{hyperref}
\hypersetup{colorlinks=true,linkcolor=black,urlcolor=black,pdfborder={0 0 0}}

% --- Métadonnées du lot -----------------------------------------------------
% Reprises telles quelles par le script de rendu, qui les affiche en tête de
% la vue de lecture. Elles fixent ce que le fichier prétend couvrir : sans
% elles, une transcription flotte sans point d'attache dans un fonds de seize
% mille feuillets.

\newcommand{\folder}[1]{\gdef\grothFolder{#1}}
\newcommand{\batch}[1]{\gdef\grothBatch{#1}}
\newcommand{\leaves}[2]{\gdef\grothFirst{#1}\gdef\grothLast{#2}}
\newcommand{\foldertitle}[1]{\gdef\grothFoldertitle{#1}}
\gdef\grothFolder{?}\gdef\grothBatch{?}\gdef\grothFirst{?}\gdef\grothLast{?}\gdef\grothFoldertitle{}

% --- Le résumé --------------------------------------------------------------
%
% Il ouvre la lecture modernisée, sous le titre « Résumé », et il est composé
% en petit. Ce n'est pas un ornement : c'est le seul passage du document qui
% ne soit pas des mathématiques, et un lecteur doit pouvoir le reconnaître
% avant de l'avoir lu — puis le sauter s'il connaît déjà le sujet, ou s'y
% arrêter s'il ne le connaît pas. Le filet qui le suit marque la reprise du
% corps.

\newenvironment{resume}
  {\small\setlength{\parskip}{0.45em}}
  {\par\medskip\hrule\medskip}

% --- Mots-clés ---------------------------------------------------------------
%
% En anglais, en clôture du résumé de la lecture modernisée : le vocabulaire
% moderne sous lequel chercher ce que les feuillets construisent. Le manifeste
% du site (`npm run manifest`) les extrait pour étiqueter la cote entière —
% cette ligne est la source unique des tags, il n'y a pas de fichier à côté.
\newcommand{\keywords}[1]{\par\smallskip\noindent
  {\footnotesize\textsc{Keywords} --- \textit{#1}}\par}

% --- L'appareil critique ----------------------------------------------------

% Le numéro de feuillet, en marge. C'est l'ancre qui permet de régler un
% désaccord sur une formule en regardant *un* feuillet plutôt qu'une chemise
% de six cents. Le site s'en sert aussi pour tourner les pages du fac-similé
% quand on fait défiler la transcription.
\newcommand{\leaf}[1]{%
  \marginnote{\footnotesize\color{gray!70}\textsf{#1}}%
  \label{leaf:#1}\ignorespaces}

% Illisible : on ne devine pas. Un mot inventé qui se lit comme les autres est
% le pire résultat possible.
\newcommand{\ill}{\textcolor{red!60!black}{[\,\ldots\,]}}

% Lecture douteuse : le mot est donné, et signalé comme incertain.
\newcommand{\uncertain}[1]{\uline{#1}}

% Ajout de l'éditeur — une lettre manquante, un mot sous-entendu.
\newcommand{\add}[1]{\textcolor{blue!50!black}{[#1]}}

% Biffé par Grothendieck lui-même. Ce qu'il a rayé fait partie du document :
% une rature dit souvent où la pensée a changé de direction.
\newcommand{\struck}[1]{\sout{#1}}

% Note du transcripteur, sur l'état matériel du feuillet.
\newcommand{\note}[1]{\par\smallskip{\small\color{gray!60!black}\textit{[#1]}}\par\smallskip}

% Note marginale de Grothendieck, distincte du corps du texte.
\newcommand{\marginal}[1]{\par\smallskip\noindent\hspace{1em}%
  {\small\color{gray!60!black}#1}\par\smallskip}

% --- Titre ------------------------------------------------------------------

\AtBeginDocument{%
  \begin{center}
    {\large\bfseries Fonds Alexandre Grothendieck --- cote n\textsuperscript{o}~\grothFolder}\\[2pt]
    {\itshape\grothFoldertitle}\\[4pt]
    {\small feuillets \grothFirst--\grothLast\ (lot \grothBatch)}\\[2pt]
    {\footnotesize\color{gray!60!black}Transcription établie d'après le fac-similé numérisé
     par l'Université de Montpellier.\\
     Les lectures incertaines sont soulignées, les illisibles notées [\,\ldots\,],
     les ajouts entre crochets.}
  \end{center}
  \vspace{1em}}

\endinput


\folder{160}
\pages{1}{22}
\dating{s.d. — le groupe « Dérivateurs » (157-1 à 160) est daté 1990-[1991]}
\watermark{Édition de démonstration\\interprétation personnelle de l'œuvre}
\foldertitle{Schémas [en] groupes : notes manuscrites (s.d.) — lecture modernisée du dossier entier}

\begin{document}

\section*{Résumé}

\begin{resume}
En caractéristique $p$, l'élévation à la puissance $p$ est additive :
$(x+y)^p = x^p + y^p$, parce que tous les coefficients binomiaux du milieu
sont divisibles par $p$. C'est une identité élémentaire, et c'est le ressort
de tout ce dossier. Elle entraîne que tout polynôme de la forme
$f(x) = c_n x^{p^n} + \dots + c_1 x^p + c_0 x$ — on les appelle les
\emph{polynômes additifs} — vérifie $f(x+y) = f(x) + f(y)$ : ce n'est pas
seulement une fonction sur la droite, c'est un endomorphisme du groupe
additif. Ses zéros forment donc un sous-groupe.

Mais « ses zéros » ne veut pas dire ici ce qu'il voudrait dire sur les
nombres réels. Prenons $f(x) = x^p$ : le seul point où il s'annule est $0$,
et pourtant il s'y annule $p$ fois. L'objet que le polynôme découpe n'est pas
l'ensemble de ses racines, c'est le \emph{schéma} qu'il définit, ici
l'algèbre $k[x]/(x^p)$, de dimension $p$ sur $k$ alors qu'elle n'a qu'un
point. On le note $\alpha_p$. Il est invisible pour qui ne regarde que les
points, et il est un vrai sous-groupe : c'est le prototype des groupes
infinitésimaux, ceux qui vivent entièrement dans un voisinage de l'origine.
À côté de lui, le polynôme $x^p - x$ découpe les $p$ éléments de
$\mathbb{F}_p$, un sous-groupe honnêtement constitué de points, et le
polynôme $x^p - a x$ pour $a$ quelconque interpole entre les deux cas. Il
faut compter non pas les points mais le \emph{rang}, c'est-à-dire la
dimension de l'algèbre : $\alpha_p$ et $\mathbb{Z}/p$ sont tous deux de
rang $p$, et pourtant ils ne se ressemblent pas.

La question du dossier est alors naturelle : les sous-groupes finis du groupe
additif $\mathbb{G}_a$ sont-ils \emph{tous} de cette forme, noyaux de
polynômes additifs ? Et si oui, comment la famille de tous les sous-groupes
d'un rang donné est-elle elle-même organisée ? La réponse, démontrée aux
pages 2 à 6, est aussi nette qu'on peut l'espérer : pour tout schéma de base
$S$ de caractéristique $p$, les sous-groupes finis localement libres de rang
$p^n$ de $\mathbb{G}_{a,S}$ sont exactement les noyaux des
$f_a(x) = x^{p^n} + a_{n-1}x^{p^{n-1}} + \dots + a_0 x$, et chacun ne
provient que d'un seul $a$. La famille de tous ces sous-groupes est donc
l'espace affine de dimension $n$, sans point double et sans point manquant.
Le sous-groupe est étale — c'est-à-dire fait de points, sans partie
infinitésimale — exactement quand $a_0$ est inversible.

Ce théorème une fois acquis, le dossier le retourne et s'en sert. Sur un
corps, se donner un sous-groupe étale de rang $p^n$ revient à se donner $n$
éléments $\xi_1, \dots, \xi_n$ linéairement indépendants sur $\mathbb{F}_p$,
et le polynôme qui s'annule exactement sur leurs combinaisons a des
coefficients que la règle de Cramer calcule. Le déterminant qui apparaît au
dénominateur — le \emph{déterminant de Moore}, formé des lignes $\xi_i$,
$\xi_i^p$, $\xi_i^{p^2}$, et ainsi de suite — se factorise complètement en
formes linéaires, et les coefficients obtenus sont des polynômes en les
$\xi_i$ invariants sous tous les changements de base à coefficients dans
$\mathbb{F}_p$. Les pages 9 à 12 montrent qu'il n'y en a pas d'autres : les
polynômes invariants sous $GL(n, \mathbb{F}_p)$ forment eux-mêmes une algèbre
de polynômes, engendrée par ces $n$ coefficients. C'est le théorème des
\emph{invariants de Dickson}, et il tombe ici comme un corollaire de la
classification des sous-groupes, par un argument de géométrie — surjectivité,
fibres finies, Main Theorem de Zariski — et non par un calcul d'algèbre.

Les pages 13 à 16 refont tout cela au-dessus d'une base quelconque, où les
$n$ éléments indépendants ne peuvent plus être écrits en ligne : ils
deviennent un faisceau localement libre $E$ de rang $n$ muni d'un
isomorphisme avec son tiré en arrière par le Frobenius. Le déterminant de
Moore devient une application polynomiale à valeurs dans la puissance
extérieure maximale, et le fait que $A_0$ soit la puissance $(p-1)$-ième du
déterminant devient une trivialisation canonique — c'est ce qui permet de
parler d'un torseur sous $\mathbb{F}_p^*$ là où on parlait d'une racine
$(p-1)$-ième. La page 16 s'arrête sur un tableau à deux colonnes qui oppose
le cas étale, entièrement réglé, au cas localement libre général ; et la
flèche qui devrait fermer la seconde colonne porte un point d'interrogation
qu'il ne lève pas.

Le reste du dossier est d'une autre nature. La page 1, isolée, énonce un
lemme d'algébrisation le long d'un fermé qui ne se rattache à rien d'autre
ici. La page 18 esquisse une descente le long d'un pincement — on recolle
deux points d'un schéma en un seul, et on demande ce qui redescend — et
s'arrête au milieu. Les pages 19 à 22, au crayon puis à l'encre, sont du
calcul à la main : le cas $n = 2$, deux racines $\alpha$ et $\beta$ dont on
élimine les coefficients, et à la dernière ligne de la page 21 le déterminant
de Moore de $(f, g, X)$ reparaît, celui-là même de la page 7. La page 22,
écrite des deux bouts du feuillet, fabrique l'exemple que la page 18
demandait : deux points où les fibres d'un sous-groupe sont abstraitement
isomorphes mais où les facteurs se croisent. Elle finit sur l'opérateur
$\lambda \mapsto (\lambda/f)^p - (\lambda/f)$, un Artin--Schreier tordu dont
le noyau est exactement la droite $\mathbb{F}_p f$ — la brique élémentaire
dont tout le reste est composé.

Un mot sur ce que ces feuillets ne disent pas. Le dossier est rangé dans un
groupe que les archivistes datent de 1990-[1991] et intitulent
« Dérivateurs » ; or rien de ce qui est écrit sur ces vingt-deux pages ne
relève des dérivateurs. C'est de la théorie des schémas en groupes finis en
caractéristique $p$, le matériau de SGA 3. C'est un constat de contenu, que
n'importe quel lecteur de la transcription peut refaire. Il ne date rien : un
mathématicien peut revenir à tout âge sur un sujet ancien, et un dossier
d'archives est une chemise physique, pas une unité thématique. L'un des deux
transcripteurs ajoute une observation d'écriture — main rapide et droite au
stylo-plume, abréviations et notation des années EGA — mais une observation
d'écriture faite sur un fac-similé est fragile, et aucun des vingt-deux
feuillets ne porte de date, d'en-tête ni de verso réutilisé qui trancherait.
La datation portée en tête de cette lecture reste donc celle de Montpellier,
inchangée, et rien ici ne prétend la corriger.

Les noms modernes sous lesquels chercher : polynômes additifs et anneau de
polynômes tordus, schémas en groupes finis en caractéristique $p$,
$\alpha_p$ et suite connexe-étale, déterminant de Moore, invariants de
Dickson, modules à Frobenius inversible, théorie d'Artin--Schreier,
pincement au sens de Ferrand.

\keywords{additive polynomial, twisted polynomial ring, finite group scheme, additive group scheme, infinitesimal group scheme, connected-etale sequence, Moore determinant, Dickson invariant, modular invariant theory, Zariski's main theorem, valuative criterion of properness, unit-root F-crystal, Artin-Schreier theory, companion matrix, Ferrand pinching, formal completion, schematic density}
\end{resume}

\pagerange{1}{22}
\section*{Le fil du dossier, et l'ordre des feuillets}

Le dossier n'est pas un texte suivi, et il ne se lit pas comme tel. Il y a
un centre, et autour de lui deux appendices et deux chantiers.

\pagerange{1}{22}
\subsection*{Les stations}

\begin{itemize}
\item \textbf{Page 1} — un lemme d'algébrisation le long d'un fermé, isolé.
Il ne se rattache à rien d'autre du dossier, et la moitié de sa prose est
perdue. On le donne pour ce qu'il est.
\item \textbf{Pages 2 à 6} — le cœur. Construction de l'application
$a \mapsto \operatorname{Ker} f_a$, et démonstration qu'elle est bijective :
surjectivité sur un corps algébriquement clos, monomorphie par déformation
aux nombres duaux, propreté par le critère valuatif, description du lieu
étale, calcul de l'espace tangent de Zariski à l'origine pour $n = 1$, puis
l'énoncé du théorème au-dessus d'une base quelconque.
\item \textbf{Pages 7 à 12} — la conséquence combinatoire. Le déterminant de
Moore et sa factorisation, la règle de Cramer qui donne les coefficients, la
divisibilité $\Delta_i = A_i \Delta_n$, et le théorème de Dickson sur les
invariants de $GL(n, \mathbb{F}_p)$, avec sa variante pour
$SL(n, \mathbb{F}_p)$.
\item \textbf{Pages 13 à 16} — la même chose au-dessus d'une base, en termes
de faisceaux. Se termine sur un tableau à deux colonnes dont la seconde
reste ouverte.
\item \textbf{Page 18} — une esquisse séparée : descente d'un sous-groupe de
$\mathbb{G}_a$ le long d'un pincement. Inachevée. La page 17 est blanche.
\item \textbf{Pages 19 à 22} — du calcul. Deux feuillets au crayon, puis
deux à l'encre, sur le cas $n = 2$ ; la page 22 rejoint la page 18.
\end{itemize}

Ce qui fait l'unité du dossier n'est pas l'ordre des feuillets mais un
objet : le polynôme additif dont les racines sont exactement un
$\mathbb{F}_p$-sous-espace donné. Il est construit abstraitement aux pages 2
à 6, écrit comme un quotient de déterminants aux pages 7 à 9, rendu
intrinsèque aux pages 13 à 16, et calculé à la main pour $n = 2$ aux
pages 19 à 22. Les quatre stations sont quatre présentations du même objet.

\pagerange{1}{22}
\subsection*{Conventions, valables pour tout le dossier}

On écrit $\mathbb{A}^n$ l'espace affine de dimension $n$ sur $\mathbb{F}_p$,
là où les feuillets écrivent $\mathbb{E}^n$, et
$\mathbb{G}_m \times \mathbb{A}^{n-1}$ l'ouvert où la première coordonnée est
inversible, là où ils écrivent $\mathbb{E}^* \times \mathbb{E}^{n-1}$.

Pour $a = (a_0, \dots, a_{n-1})$ à coefficients dans un anneau $R$ de
caractéristique $p$, on pose
\[
  f_a(x) = x^{p^n} + a_{n-1}x^{p^{n-1}} + \dots + a_1 x^p + a_0 x ,
\]
endomorphisme du groupe $\mathbb{G}_{a,R}$, fini localement libre de degré
$p^n$ puisque le polynôme est unitaire de degré $p^n$. Son noyau
$\operatorname{Ker} f_a$ est donc un sous-schéma en groupes fini localement
libre de rang $p^n$.\footnote{Le langage moderne le dit d'un mot :
$\operatorname{End}(\mathbb{G}_{a,R})$ est l'anneau de polynômes tordus
$R\{F\}$, avec $F c = c^p F$, et $f_a$ en est un élément unitaire de
$F$-degré $n$. Les feuillets ne nomment pas cet anneau mais en manipulent
constamment la loi de composition.}

$F$ désigne le Frobenius, $\alpha_{p^r} = \operatorname{Ker}(F^r)$ le noyau
de sa $r$-ième itérée, de rang $p^r$ et à un seul point. Il ne faut pas le
confondre avec $\alpha_p^r$, qui est un produit de $r$ copies de $\alpha_p$
et n'est pas isomorphe à $\alpha_{p^r}$ dès que $r \geqslant 2$ ; le dossier
utilise partout le premier.

$X_n$ désigne, à partir de la page 2 et jusqu'à la fin, le schéma des
sous-groupes finis localement libres de rang $p^n$ de $\mathbb{G}_a$. Mais
sur la page 1, et seulement là, la même lettre indicée désigne les
voisinages infinitésimaux d'un fermé. Les deux objets n'ont rien de commun,
et pour éviter qu'on lise le second comme le premier on note ici $X^{(m)}$
le $m$-ième voisinage infinitésimal de la page 1.

Enfin, on suit la convention de EGA : pour $E$ localement libre,
$V(E) = \operatorname{Spec} \underline{\operatorname{Sym}}(E)$, de sorte
qu'une section de $V(E)$ au-dessus de $S$ est un homomorphisme
$E \to \mathcal{O}_S$, c'est-à-dire une section du \emph{dual}
$\check{E}$. C'est aussi ce que font les pages 13 et 14, qui parlent de
sections de $\check{E}$ et de points de $V(E)$ sans jamais les
distinguer.\footnote{La variance est ici tout sauf un détail : les
opérations $\xi \mapsto \xi^{(p)}$ des pages 14 et 15 vivent sur $\check{E}$
et utilisent la structure $\check{u} : \check{E}^{(p)} \to \check{E}$
déduite de $u : E^{(p)} \xrightarrow{\ \sim\ } E$ par dualité. Les lire sur
$E$ donnerait des énoncés faux.}

\pagerange{1}{22}
\subsection*{Ce que le contenu montre, et ce qu'il ne date pas}

Il faut le dire ici plutôt que de le taire. Le dossier appartient, dans
l'inventaire de Montpellier, à un groupe intitulé « Dérivateurs » et daté
1990-[1991]. Rien sur ces vingt-deux pages ne relève des dérivateurs, ni des
catégories dérivées, ni de l'algèbre homotopique. Ce qui s'y trouve est de
la théorie des schémas en groupes finis en caractéristique $p$ : polynômes
additifs, $\alpha_p$, Verschiebung, critère valuatif, Main Theorem de
Zariski, invariants de Dickson, complétion formelle le long d'un fermé. Le
lemme de la page 1 est écrit dans le vocabulaire de EGA III et IV ; les
pages 2 à 16 sont, mot pour mot, le genre de matériau que SGA 3 organise.

Ce constat est solide en ceci qu'il est vérifiable : il suffit de lire la
transcription. Il faut être exact sur sa portée, qui est étroite. Il dit ce
que les feuillets traitent ; il ne dit pas quand ils ont été écrits. Un
mathématicien revient sur un sujet ancien quand il lui plaît, et un dossier
d'archives est une chemise de papier, groupée par les archivistes selon la
place qu'elle occupait, non selon ce qu'elle contient. Un contenu de 1960 ne
prouve pas une écriture de 1960.

À côté de ce constat, l'un des deux transcripteurs porte une observation
d'une autre espèce : une main rapide et droite, au stylo-plume, avec les
abréviations et la notation des années EGA. C'est une observation
paléographique, et elle est fragile — elle a été faite sur un fac-similé
numérique, où l'encre, la pression, le papier et le format ne se jugent
pas. Elle demanderait l'œil d'un humain sur les feuillets eux-mêmes. Le même
transcripteur note qu'aucune des vingt-deux pages ne porte de date, d'en-tête
ni de verso réutilisé qui trancherait la question.

En conséquence : cette lecture ne conclut aucune date. La ligne de datation
portée en tête est celle des archivistes, reprise verbatim des deux
transcriptions, et elle y reste. Ce qui est offert ici est la description de
ce que les feuillets montrent, pour que quelqu'un qui pourra les manipuler
sache quoi regarder.

\pagerange{1}{1}
\section*{Un lemme d'algébrisation isolé (page 1)}

Le feuillet est seul de son espèce dans le dossier, et il est très abîmé :
la moitié de sa prose manque, la démonstration s'arrête au milieu d'une
phrase, et une parenthèse ouverte au point 2) n'est jamais fermée. On donne
ce qui se tient.

Soient $f : X' \to X$ un morphisme de type fini entre schémas
noethériens,\footnote{Le mot « propre » est écrit puis biffé sur le
feuillet, et remplacé par « de type fini ». On suit la correction.} $V'
\subset X'$ l'ouvert des points où $f$ est étale, $Z \subset X$ une partie
fermée, $\hat{X}$ le complété formel de $X$ le long de $Z$, et $X^{(m)}$ le
$m$-ième voisinage infinitésimal de $Z$ dans $X$. On se donne une section
formelle $q = (q_m)$ de $X'$ au-dessus de $\hat{X}$, c'est-à-dire une famille
compatible de sections de $X' \times_X X^{(m)}$ au-dessus de $X^{(m)}$, et on
note $Z'$ l'image de la section au niveau réduit, de sorte que $f$ induise
$Z' \xrightarrow{\ \sim\ } Z$.

La page pose alors une question de densité schématique et énumère trois
conditions dont elle dit qu'elles sont liées :

\begin{itemize}
\item $q_m^{-1}(V')$ est schématiquement dense dans $X^{(m)}$ ;
\item $X'$ est l'adhérence schématique de la famille des $q_m(X^{(m)})$ ;
\item la famille des $V'$ vus dans les $X^{(m)}$ est schématiquement dense
dans $X'$, ce que la marge résume en « i.e. en un $\hat{V'}$ ».
\end{itemize}

Le cas retenu est celui où $Z \subset V'$, c'est-à-dire où $f$ est étale aux
points de $Z$ : il existe alors un ouvert $U' \subset V'$ contenant $Z'$ tel
que $U' \times_X Z \xrightarrow{\ \sim\ } Z' \xrightarrow{\ \sim\ } Z$. La
démonstration esquissée scinde la troisième condition en deux, une de
densité schématique de $V'$ dans $X'$ et une de densité de $\hat{V'}$ dans
$U'$, et s'interrompt là.\footnote{Ce qui manque n'est pas récupérable :
les mots qui relient les trois conditions sont illisibles, et l'implication
que la page veut établir — laquelle des trois entraîne laquelle — n'est
lisible sur aucune ligne. On ne la reconstruit pas. Le contexte général est
celui des théorèmes d'algébrisation le long d'un fermé, autour de EGA III 5
et des voisinages étales ; la page n'en nomme aucun.}

\pagerange{2}{6}
\section*{Les sous-groupes finis de $\mathbb{G}_a$ et l'application $\varphi_n$ (pages 2 à 6)}

Le but est l'application
\[
  \varphi_n : \mathbb{A}^n \longrightarrow X_n, \qquad
  a = (a_0, \dots, a_{n-1}) \longmapsto \operatorname{Ker} f_a ,
\]
et l'affirmation qu'elle est un isomorphisme. Les pages 2 à 5 la démontrent
par morceaux : elle est surjective, elle est un monomorphisme, elle est
propre ; d'où un isomorphisme sur $X_{n,\mathrm{réd}}$. La page 5 pose alors
la question de la réduction de $X_n$, et la page 6 énonce le théorème
complet.\footnote{Le dossier utilise $X_n$ sans jamais le construire. Il est
représentable : un sous-schéma fermé fini localement libre de rang $d$ de
$\mathbb{A}^1_S$ est le lieu des zéros d'un unique polynôme unitaire de
degré $d$ — le polynôme caractéristique de la multiplication par $x$ sur son
algèbre —, ce qui plonge la famille dans l'espace affine $\mathbb{A}^{p^n}$
des polynômes unitaires de degré $p^n$ ; la condition d'être un sous-groupe y
découpe $X_n$. C'est cette représentabilité qui donne un sens aux mots
« propre », « quasi-fini » et « réduit » employés plus bas.}

\pagerange{2}{2}
\subsection*{Surjectivité (page 2)}

Il suffit de voir que $\mathbb{A}^n(k) \to X_n(k)$ est surjective pour $k$
algébriquement clos. Soit donc $H \subset \mathbb{G}_{a,k}$ fini de rang
$p^n$. Sa composante neutre $H^0$ est un sous-schéma fini connexe de
$\mathbb{A}^1_k$ concentré à l'origine, donc de la forme
$\operatorname{Spec} k[x]/(x^{p^r})$, c'est-à-dire
$\alpha_{p^r} = \operatorname{Ker}(F^r)$ ; le quotient $H/H^0$ est étale, tué
par $p$, donc $(\mathbb{Z}/p)^s$ avec $r + s = n$, et sur $k$ parfait la
suite connexe-étale est scindée :
\[
  H \;\simeq\; \alpha_{p^r} \times (\mathbb{Z}/p)^s_k , \qquad r + s = n .
\]
\footnote{La transcription lit $\alpha_p^r$ à cet endroit. C'est
$\alpha_{p^r}$ qu'il faut, et c'est ce que la page écrit quatre lignes plus
bas, sous la forme $\alpha_{p^r}$ glosée
« $\operatorname{Ker} \mathrm{Frob}^r$ ». Les deux groupes sont différents
dès $r \geqslant 2$ — $\alpha_{p^2}$ a une algèbre $k[x]/(x^{p^2})$, $\alpha_p^2$
une algèbre $k[x,y]/(x^p, y^p)$ —, et seul $\alpha_{p^r}$ se plonge dans
$\mathbb{G}_a$. La distinction entre indice et exposant se joue, dans une
main rapide, sur la hauteur d'un caractère.}

Soient $\xi_1, \dots, \xi_s$ une base de $H(k)$ sur $\mathbb{F}_p$. Le
sous-groupe constant qu'ils engendrent est le lieu des zéros du polynôme
\[
  f_\alpha(x) = \prod_{\lambda \in \mathbb{F}_p^s}
  \Bigl(x - \sum_{i=1}^{s} \lambda_i \xi_i\Bigr),
\]
qui est additif précisément parce que son ensemble de racines est un
$\mathbb{F}_p$-sous-espace, et séparable de degré
$p^s$.\footnote{La page écrit « $H(k)_k = \operatorname{Ker}(f_\alpha)$,
$\alpha \in \mathbb{E}^s_k$, à justifier ». La justification est celle qu'on
vient de donner, et c'est un fait classique sur les polynômes additifs
(Ore) : un polynôme unitaire séparable est additif si et seulement si son
lieu de zéros est un $\mathbb{F}_p$-sous-espace.} On pose alors
\[
  f = F^r \circ f_\alpha , \qquad
  f(x) = x^{p^{r+s}} + \alpha_{s-1}^{p^r} x^{p^{r+s-1}} + \dots
  + \alpha_0^{p^r} x^{p^r} ,
\]
d'où
\[
  \operatorname{Ker} f = f_\alpha^{-1}(\alpha_{p^r}) \;\supset\;
  H(k)_k \cdot \alpha_{p^r} = H ,
\]
et les deux membres étant de rang $p^{r+s} = p^n$, l'inclusion est une
égalité. Donc $H = \operatorname{Ker} f$ avec $f$ additif unitaire de degré
$p^n$, c'est-à-dire $H = \varphi_n(a)$ pour un $a \in
k^n$.\footnote{La page écrit $f = f_\alpha \circ \mathrm{Frob}^r$. C'est
l'autre ordre qu'il faut : avec $f = f_\alpha \circ F^r$ on aurait
$\operatorname{Ker} f = F^{-r}(\operatorname{Ker} f_\alpha)$, et non
$f_\alpha^{-1}(\alpha_{p^r})$ comme la page l'écrit deux lignes plus bas ; et
les $\xi_i$ ne seraient pas annulés. Avec $f = F^r \circ f_\alpha$ tout
s'ajuste, y compris le développement que la page donne — à ceci près que les
coefficients y sont écrits $\alpha_j$ au lieu de $\alpha_j^{p^r}$, ce qui est
sans conséquence sur $k$ parfait. C'est une inversion d'écriture, pas une
erreur de raisonnement.}

\pagerange{2}{3}
\subsection*{Monomorphie, par déformation (pages 2 et 3)}

Soient $k$ une $\mathbb{F}_p$-algèbre et $a, b$ avec
$\operatorname{Ker} f_a = \operatorname{Ker} f_b$. En notant $H$ ce noyau
commun, $f_a$ et $f_b$ se factorisent tous deux par $\mathbb{G}_{a,k}/H$, et
comme ils y induisent des isomorphismes sur leur image on obtient
$f_b = u \circ f_a$ avec $u \in \operatorname{Aut}(\mathbb{G}_{a,k})$. Tout
revient donc à :

\textbf{Lemme.} Si $u \in \operatorname{Aut}(\mathbb{G}_{a,k})$ est tel que
$u \circ f_a$ soit encore de la forme $f_b$, alors
$u = \operatorname{id}$.

Si $k$ est réduit, c'est immédiat : $\operatorname{Aut}(\mathbb{G}_{a,k})$
se réduit alors aux homothéties $k^*$, et
$u \circ f_a = c_0 f_a$ ne peut être unitaire que si $c_0 = 1$.\footnote{En
termes de l'anneau de polynômes tordus : $R\{F\}^\times = R^\times$ dès que
$R$ est réduit, parce que le $F$-degré est alors additif. C'est exactement ce
qui tombe en défaut sur un anneau non réduit, et c'est pourquoi la page passe
aux déformations.}

Le cas général se ramène, $k$ étant supposé noethérien, à une suite
d'extensions de carré nul.\footnote{La page écrit « OPS $k$ noethérien » —
on peut supposer — et ne détaille pas la réduction. Elle est standard :
l'énoncé est local et de présentation finie, donc se descend à une
$\mathbb{F}_p$-algèbre de type fini, où le nilradical est nilpotent et se
dévisse en idéaux de carré nul.} Soit donc $J \subset k$ un idéal avec
$J^2 = 0$, $k_0 = k/J$, et supposons $u_0 = \operatorname{id}$ sur
$\mathbb{G}_{a,k_0}$. Alors
\[
  u(x) = c_0 x + c_1 x^p + c_2 x^{p^2} + \dots + c_N x^{p^N},
  \qquad c_0 \in 1 + J, \quad c_i \in J \ (i \geqslant 1) ,
\]
et en développant
\[
  u \circ f_a(x) = \sum_{j=0}^{N} c_j
  \Bigl(x^{p^{n+j}} + a_{n-1}^{p^j} x^{p^{n+j-1}} + \dots
  + a_0^{p^j} x^{p^j}\Bigr).
\]
Le coefficient de $x^{p^{n+N}}$ vaut $c_N$ ; comme $u \circ f_a = f_b$ est de
degré $p^n$, il est nul, donc $c_N = 0$. Le coefficient de $x^{p^{n+N-1}}$
vaut alors $c_{N-1}$, nul de même, et de proche en proche
$c_1 = \dots = c_N = 0$. Il reste $u(x) = c_0 x$, et $c_0 = 1$ parce que
$f_b$ est unitaire.\footnote{La page écrit à cet endroit
« $u(x) = c_0(x) = c_0 x^{p^n}$ » — l'exposant est de trop, et la
conclusion voulue est bien $u(x) = c_0 x$ puis $c_0 = 1$. La ligne suivante
du feuillet, « et il faut aussi $c_0 = 1$ », le confirme.}

\pagerange{3}{4}
\subsection*{Propreté, par le critère valuatif (pages 3 et 4)}

\textbf{Lemme.} Soient $S = \operatorname{Spec} V$ un trait sur
$\mathbb{F}_p$, de point générique $\eta$ et de corps des fractions $K$, et
$H \subset \mathbb{G}_{a,S}$ fini plat de rang $p^n$ sur $S$. Si
$H_\eta = \operatorname{Ker} f_{a_\eta}$ avec $a_\eta \in K^n$, alors
$a_\eta \in V^n$, et $H = \operatorname{Ker} f_{a_\eta}$ sur $S$ tout entier.

\emph{Démonstration.} Supposons $a_\eta \notin V^n$ et posons
$\nu = -\inf_i \operatorname{val}(a_i) > 0$, $\pi$ une uniformisante,
$b_i = \pi^\nu a_i$ : tous les $b_i$ sont dans $V$ et l'un d'eux est une
unité. Soit
\[
  f(x) = \pi^\nu x^{p^n} + b_{n-1}x^{p^{n-1}} + \dots + b_1 x^p + b_0 x ,
  \qquad f : \mathbb{G}_{a,S} \to \mathbb{G}_{a,S} .
\]
Sur la fibre générique $f_\eta = \pi^\nu f_{a_\eta}$ et $\pi^\nu$ est
inversible, donc $\operatorname{Ker} f_\eta = H_\eta$. Le morphisme $f$ est
plat, parce qu'il l'est fibre à fibre : sur la fibre spéciale $f_s$ n'est pas
nul, un des $b_{i,s}$ étant non nul. Par suite
$\operatorname{Ker} f = f^{-1}(0)$ est plat sur $S$, de fibre générique
$H_\eta$ ; c'est donc l'adhérence schématique de $H_\eta$ dans
$\mathbb{G}_{a,S}$, c'est-à-dire $H$ lui-même. La fibre spéciale
$\operatorname{Ker} f_s$ devrait donc être de rang $p^n$. Or $\pi^\nu$
s'annule sur la fibre spéciale, de sorte que
\[
  f_s(x) = b_{m,s} x^{p^m} + \dots , \qquad
  m = \max\{\, i \leqslant n-1 \;:\; b_{i,s} \neq 0 \,\} ,
\]
et $\operatorname{Ker} f_s$ est de rang $p^m$ avec $m \leqslant n - 1$.
Contradiction.\footnote{La page conclut « il est de rang $p^{n-i}$ si
$b_{i,s}$ est le premier des $b_{0,s}, b_{1,s}, \dots$ qui soit non nul ».
Le rang du noyau d'un polynôme additif est son degré, gouverné par le
\emph{dernier} indice non nul et non par le premier ; et c'est bien
$p^m$, $m$ maximal, qu'il faut écrire. La conclusion est la même — le rang
tombe strictement en dessous de $p^n$ dès que $\pi^\nu$ meurt — mais l'énoncé
intermédiaire de la page n'est pas celui-là.}

\emph{Une fois ce lemme acquis}, $\varphi_n$ est propre par le critère
valuatif, et un monomorphisme propre est une immersion fermée. Comme elle est
de plus surjective, elle induit
\[
  \mathbb{A}^n \xrightarrow{\ \sim\ } X_{n,\mathrm{réd}} .
\]

\pagerange{4}{5}
\subsection*{Le lieu étale (pages 4 et 5)}

Soit $X_n^{\text{ét}} \subset X_n$ l'ouvert des sous-groupes étales sur la
base. Alors $\varphi_n$ induit
\[
  \mathbb{G}_m \times \mathbb{A}^{n-1} \xrightarrow{\ \sim\ }
  X_n^{\text{ét}} ,
\]
l'ouvert de gauche étant celui des $(a_0, \dots, a_{n-1})$ à $a_0$
inversible. La raison est d'une ligne : $f_a'(x) = a_0$, donc $f_a$ est étale
— et son noyau aussi — exactement quand $a_0$ est
inversible.\footnote{La page porte « à justifier » en marge de l'existence de
l'ouvert $X_n^{\text{ét}}$, non de l'énoncé. Que le lieu où un schéma en
groupes fini localement libre est étale sur la base soit ouvert est un fait
général ; c'est le lieu où le discriminant est inversible.}

\pagerange{5}{6}
\subsection*{$X_n$ est-il réduit ? L'espace tangent à l'origine pour $n = 1$ (pages 5 et 6)}

La page pose la question sous cette forme : $X_n$ est-il réduit, c'est-à-dire
$\varphi_n$ est-elle une immersion fermée dans les deux sens ?\footnote{La
page ajoute, sous forme de seconde question, « ou $X_n^{\mathrm{réd}}
\simeq \mathbb{E}^* \times \mathbb{E}^{n-1}$ ? ». Ce ne peut être
$X_n^{\mathrm{réd}}$ — la page vient d'établir que celui-ci est $\mathbb{A}^n$
tout entier — mais $X_n^{\text{ét}}$, dont le lemme précédent donne
précisément cet isomorphisme. On lit l'exposant comme une reprise distraite
de la ligne du dessus.}

Elle traite $n = 1$ au point $0$, ce qui revient à montrer que l'espace
tangent de Zariski de $X_1$ en $0$ est de dimension $1$ — il est a priori de
dimension au moins $1$. Il faut donc décrire les points de $X_1$ à valeurs
dans les nombres duaux $k = \mathbb{F}_p[J]/J^2$ qui se réduisent sur $0$,
c'est-à-dire les sous-groupes plats $H \subset \mathbb{G}_{a,k}$ tels que
$H_0 = \alpha_p$.

Le calcul de la page procède ainsi. Posons $G = \mathbb{G}_{a,k}$,
$G_0 = \mathbb{G}_{a,\mathbb{F}_p}$, $K = G/H$ et $K_0 = G_0/H_0$. Le
quotient $K$ est une extension de $K_0$ par $J K_0$, et l'obstruction à la
scinder — donc la donnée de $H$ au-dessus de $H_0$ — est mesurée par un
homomorphisme
\[
  u : H_0 \longrightarrow J K_0 .
\]
Ici $H_0 = \alpha_p \subset G_0 = \mathbb{G}_a$ et
$K_0 \simeq J K_0 \simeq \mathbb{G}_a$, de sorte que l'extension est triviale
et que l'espace tangent cherché est
\[
  T_0 X_1 \;=\; \operatorname{Hom}_{\mathbb{F}_p\text{-gr}}(\alpha_p,
  \mathbb{G}_a) .
\]
Or tout homomorphisme de $\alpha_p$ dans $\mathbb{G}_a$ est tué par le
Frobenius au but, donc se factorise par $\alpha_p$, et
$\operatorname{End}(\alpha_p) = \mathbb{F}_p \cdot \operatorname{id}$ sur
$\mathbb{F}_p$ : l'espace tangent est de dimension $1$, et $X_1$ est lisse
en $0$.\footnote{Cette dernière étape n'est pas sur la page, qui s'arrête sur
l'expression $\operatorname{Hom}_{\mathbb{F}_p}(H_0, J K_0)$ sans en calculer
la dimension. Elle est nécessaire pour que le calcul réponde à la question
posée, et elle est courte ; on la donne. Les six lignes de la page 5 qui
mènent à $u$ sont, elles, parmi les plus abîmées du dossier : trois mots sur
quatre sont illisibles, et on n'a restitué que le squelette de l'argument —
extension, scindage, obstruction — que les formules encore lisibles
soutiennent. La page 19 reprend le même calcul au crayon, et ne l'éclaircit
pas.}

\pagerange{6}{6}
\subsection*{Le théorème (page 6)}

\textbf{Théorème.} Soit $S$ un schéma de caractéristique $p$. Pour tout
sous-schéma en groupes $H \subset \mathbb{G}_{a,S}$ fini localement libre de
rang $p^n$, il existe un unique
$a = (a_0, \dots, a_{n-1}) \in \Gamma(S, \mathcal{O}_S)^n$ tel que
$H = \operatorname{Ker} f_a$. Autrement dit
$\varphi_n : \mathbb{A}^n \xrightarrow{\ \sim\ } X_n$. De plus $H$ est étale
sur $S$ si et seulement si $a_0 \in \Gamma(S, \mathcal{O}_S^*)$.

L'énoncé est vrai.\footnote{La démonstration se complète ainsi : un
sous-schéma fermé de $\mathbb{A}^1_S$ fini localement libre de rang $p^n$ est
le lieu des zéros d'un unique polynôme unitaire $f$ de degré $p^n$ (polynôme
caractéristique de la multiplication par $x$), et la condition que ce lieu
soit un sous-groupe force $f$ additif. Le dossier n'emprunte pas ce chemin :
il démontre l'isomorphisme sur le réduit, pose la question de la réduction de
$X_n$, et y répond au seul cas $n = 1$ et au seul point $0$. L'énoncé général
est donc, dans ce dossier, annoncé et non établi ; c'est la plus grande
lacune du feuillet, et elle ne tient pas à ce que la page soit abîmée.}

\pagerange{7}{12}
\section*{Le déterminant de Moore et les invariants de Dickson (pages 7 à 12)}

Le théorème une fois là, la question devient : \emph{quel} $a$ ? Sur un
corps, un sous-groupe étale de rang $p^n$ est le $\mathbb{F}_p$-espace
engendré par $n$ éléments indépendants, et il s'agit de calculer $a$ à
partir d'eux. La réponse fait apparaître un déterminant, et le déterminant
fait apparaître, sans qu'on l'ait cherchée, toute la théorie des invariants
du groupe linéaire sur $\mathbb{F}_p$.

\pagerange{7}{7}
\subsection*{La factorisation du déterminant de Moore (page 7)}

\textbf{Lemme.} Dans $\mathbb{F}_p[X_1, \dots, X_n]$,
\[
  \Delta_n(X_1, \dots, X_n) \;=\;
  \det\bigl(X_j^{p^{i-1}}\bigr)_{1 \leqslant i, j \leqslant n}
  \;=\; \prod_{j=0}^{n-1} \ \prod_{\alpha \in \mathbb{F}_p^{\,j}}
  \Bigl(X_{j+1} + \sum_{i=1}^{j} \alpha_i X_i\Bigr).
\]

L'argument de la page est de divisibilité et de degré. Chaque facteur
$P_\lambda = X_{j+1} + \sum \alpha_i X_i$ est irréductible, ils sont deux à
deux non associés, et chacun annule le déterminant — deux colonnes
deviennent liées —, donc le produit $P$ divise le déterminant $Q$. Les degrés
coïncident :
\[
  \deg P \;=\; 1 + p + \dots + p^{n-1} \;=\; \frac{p^n - 1}{p - 1}
  \;=\; \operatorname{card} \mathbb{P}^{n-1}(\mathbb{F}_p) ,
\]
donc $Q = c P$ avec $c \in \mathbb{F}_p^*$. La page s'arrête sur « reste à
déterminer $c$ » ; avec la normalisation ci-dessus, où chaque facteur est
unitaire en sa variable de plus grand indice, on a $c =
1$.\footnote{Vérification directe pour $n = 2$ :
$\det\begin{pmatrix} X_1 & X_2 \\ X_1^p & X_2^p \end{pmatrix}
= X_1 X_2^p - X_1^p X_2$, et le produit vaut
$X_1 \prod_{\alpha \in \mathbb{F}_p}(X_2 + \alpha X_1)
= X_1 (X_2^p - X_1^{p-1} X_2)$, les deux étant égaux. Le cas général suit
par récurrence sur $n$ en comparant les coefficients de
$X_n^{p^{n-1}}$.}\footnote{La page écrit
$\frac{p^n - 1}{p-1} = \operatorname{card} \mathbb{P}^n(\mathbb{F}_p)$. C'est
le cardinal de $\mathbb{P}^{n-1}(\mathbb{F}_p)$ que la formule donne : les
formes linéaires non nulles en $n$ variables, modulo $\mathbb{F}_p^*$, sont
les points de l'espace projectif de dimension $n-1$. Le même décalage
reparaît sur la page 20, au crayon, sous la forme
$\prod_{\mathbb{P}^n(\mathbb{F}_p)}(\lambda_1 X_1 + \dots + \lambda_n X_n)$ ;
il est donc constant chez lui dans ce dossier, et non accidentel.}

\textbf{Corollaire.} Pour $k$ un corps de caractéristique $p$ et
$\xi \in k^n$, on a $\Delta_n(\xi) = 0$ si et seulement si
$\xi_1, \dots, \xi_n$ sont linéairement dépendants sur $\mathbb{F}_p$. Par
suite l'ouvert $U_n \subset \mathbb{A}^n$ où $GL(n, \mathbb{F}_p)$ opère
librement est $U_n = \mathbb{A}^n_{\Delta_n}$.

\pagerange{8}{8}
\subsection*{Cramer (page 8)}

Cherchons $a$ tel que $f_a(\xi_i) = 0$ pour $1 \leqslant i \leqslant n$.
C'est un système de $n$ équations linéaires en $a_0, \dots, a_{n-1}$, de
seconds membres $-\xi_i^{p^n}$ et de déterminant $\pm\Delta_n(\xi)$. Si
$\Delta_n(\xi)$ est inversible — c'est-à-dire, par le corollaire précédent,
si les $\xi_i$ engendrent un sous-groupe étale de rang $p^n$ — il y a une
solution et une seule, et Cramer la donne comme un quotient de deux
déterminants : celui du bas est $\Delta_n$, celui du haut est $\Delta_j$,
obtenu en remplaçant dans la matrice de Moore la ligne des $\xi_i^{p^j}$ par
la ligne des $\xi_i^{p^n}$.

La façon propre de fixer les signes une fois pour toutes est de poser
\[
  D(T) \;=\; \det\begin{pmatrix}
    X_1 & \cdots & X_n & T \\
    X_1^p & \cdots & X_n^p & T^p \\
    \vdots & & \vdots & \vdots \\
    X_1^{p^n} & \cdots & X_n^{p^n} & T^{p^n}
  \end{pmatrix}
  \;=\; \sum_{j=0}^{n} \Delta_j(X)\, T^{p^j} ,
\]
le développement se faisant selon la dernière colonne, de sorte que le
coefficient de $T^{p^n}$ soit exactement le déterminant de Moore
$\Delta_n(X)$ et que $\Delta_j(X)$ soit, au signe $(-1)^{n+j}$ près, le
mineur obtenu en supprimant la ligne des $X_i^{p^j}$.\footnote{Cette
normalisation est la nôtre ; la page écrit les mineurs un à un, et plusieurs
de leurs lignes sont illisibles, la transcription les restituant selon la
règle de Cramer que la page énonce. Elle a l'avantage de rendre l'identité
suivante exacte, signes compris, là où la page laisse flotter des $\pm$.}

\pagerange{9}{9}
\subsection*{Les $A_i$, et pourquoi ce sont des polynômes (page 9)}

Prenons $k = \mathbb{F}_p[X_1, \dots, X_n]$ et $\xi_i = X_i$ : puisque
$D(X_i) = 0$ (deux colonnes égales),
\[
  \Delta_n X_i^{p^n} + \Delta_{n-1} X_i^{p^{n-1}} + \dots + \Delta_0 X_i = 0
  \qquad (1 \leqslant i \leqslant n) .
\]
La page affirme ensuite que les $\Delta_i$ sont divisibles par $\Delta_n$ et
pose $\Delta_i = A_i \Delta_n$. La raison, qu'elle ne donne pas, est
lisible sur la factorisation du lemme de la page 7 appliquée à $n+1$
variables :
\[
  D(T) \;=\; \Delta_n(X) \prod_{\lambda \in \mathbb{F}_p^n}
  \Bigl(T - \sum_{i=1}^n \lambda_i X_i\Bigr) ,
\]
de sorte que $D(T)/\Delta_n$ est un polynôme unitaire en $T$, additif, de
degré $p^n$, dont les coefficients sont les $A_i$.\footnote{Cette écriture
rend tout le reste immédiat, et c'est la raison de la donner : les $A_i$ sont
des polynômes, ils sont homogènes de degré $p^n - p^i$, et ils sont
invariants sous $GL(n, \mathbb{F}_p)$ parce que le produit de droite est
invariant — le groupe ne fait que permuter les $\lambda$. La page démontre
l'invariance autrement, en marge, par $\Delta_n(u \cdot X) = \det(u)
\Delta_n(X)$, ce qui marche aussi puisque les $\Delta_j$ se transforment tous
par le même facteur.}

En particulier, pour $j = 0$, le mineur est $\det(X_j^{p^i})_{1 \leqslant i
\leqslant n}$, c'est-à-dire $\Delta_n^p$, d'où
\[
  \Delta_0 \;=\; (-1)^n\, \Delta_n^{\,p} ,
  \qquad\text{donc}\qquad
  A_0 \;=\; (-1)^n\, \Delta_n^{\,p-1} .
\]
\footnote{La page écrit $\Delta_0 = \pm \Delta_n^p (-1)^n$, et le corollaire
juste en dessous $A_0 = \pm \Delta_n^{p-1}$. Les deux ne sont pas en conflit :
ils se déduisent l'un de l'autre par $A_0 = \Delta_0/\Delta_n$, et le
transcripteur l'avait déjà noté. Ce que la normalisation ci-dessus apporte
est de remplacer le $\pm$ flottant par le signe exact $(-1)^n$, qui vaut
dans les deux identités.}

\textbf{Corollaire.} $X_i^{p^n} + A_{n-1}X_i^{p^{n-1}} + \dots + A_0 X_i = 0$
pour $1 \leqslant i \leqslant n$, et les $A_i$ sont caractérisés par cette
relation.

\textbf{Corollaire.} Sur un corps $k$ de caractéristique $p$, on a
$\Delta_n(\xi) \neq 0$ — c'est-à-dire $\xi_1, \dots, \xi_n$ linéairement
indépendants sur $\mathbb{F}_p$ — si et seulement si $\Delta_0(\xi) \neq 0$.
C'est immédiat sur $\Delta_0 = (-1)^n \Delta_n^p$.

\pagerange{9}{11}
\subsection*{Le théorème de Dickson (pages 9 à 11)}

Soit $\Gamma_n = GL(n, \mathbb{F}_p)$, opérant sur $\mathbb{A}^n$ par
changement de base à coefficients dans $\mathbb{F}_p$, et
$Q_n = \mathbb{A}^n/\Gamma_n = \operatorname{Spec} A_n$ avec
$A_n = \mathbb{F}_p[X_1, \dots, X_n]^{\Gamma_n}$. Les $A_i$ étant invariants,
le morphisme
\[
  A_* : \mathbb{A}^n \longrightarrow \mathbb{A}^n, \qquad
  \xi \longmapsto (A_0(\xi), \dots, A_{n-1}(\xi))
\]
se factorise par $Q_n$.

\textbf{Théorème.} $Q_n \to \mathbb{A}^n$ est un isomorphisme,
c'est-à-dire
\[
  \mathbb{F}_p[X_1, \dots, X_n]^{GL(n, \mathbb{F}_p)}
  \;=\; \mathbb{F}_p[A_0, \dots, A_{n-1}] ,
\]
algèbre de polynômes sur les $n$ générateurs $A_i$, homogènes de degrés
$p^n - p^i$.\footnote{Ce sont les \emph{invariants de Dickson}, établis par
L. E. Dickson en 1911. La page ne le nomme pas et ne renvoie à rien. La
démonstration donnée ici est géométrique et tire tout de la classification
des sous-groupes ; ce n'est pas celle de Dickson, qui est algébrique.}

\emph{Démonstration.} $A_*$ envoie $U_n = \mathbb{A}^n_{\Delta_n}$ sur
l'ouvert $\mathbb{G}_m \times \mathbb{A}^{n-1}$ — c'est le corollaire
$A_0 = (-1)^n \Delta_n^{p-1}$ — et y est bijectif sur les orbites, donc
$Q_n \to \mathbb{A}^n$ est birationnel. $Q_n$ est intègre, $\mathbb{A}^n$ est
normal ; il suffit donc, par le Main Theorem de Zariski, de voir que
$Q_n \to \mathbb{A}^n$ est quasi-fini et surjectif, et pour cela que sur un
corps algébriquement clos $k$ l'application $A_* : k^n \to k^n$ est
surjective à fibres finies.

\emph{Fibres finies.} Les fibres sont exactement les orbites de $\Gamma_n$,
qui sont finies. La page en tire au passage, indépendamment du Main Theorem,
que $Q_n \to \mathbb{A}^n$ est surjectif, radiciel et birationnel.

\emph{Surjectivité.} Soit $a \in k^n$. Si $a = 0$, c'est l'image de $0$.
Sinon soit $r$ le premier indice avec $a_r \neq 0$. Le cas $r = 0$ est celui
du sous-groupe étale : $\operatorname{Ker} f_a$ est alors étale de rang
$p^n$, et sur $k$ algébriquement clos ses points forment un
$\mathbb{F}_p$-espace de dimension $n$ ; une base $\xi$ de cet espace vérifie
$A_j(\xi) = a_j$ pour tout $j$, puisque $D_\xi(T)/\Delta_n(\xi)$ est le
polynôme unitaire additif de lieu de zéros $\mathbb{F}_p\xi$, c'est-à-dire
$f_a$. Le cas général s'y ramène par le

\textbf{Lemme.}
\[
  A_0(X_1, \dots, X_{n-1}, 0) = 0, \qquad
  A_i(X_1, \dots, X_{n-1}, 0) = A_{i-1}(X_1, \dots, X_{n-1})^{\,p}
  \quad (1 \leqslant i \leqslant n) ,
\]
et plus généralement
\[
  A_\alpha(X_1, \dots, X_{n-s}, 0, \dots, 0) =
  \begin{cases}
    0 & \text{si } 0 \leqslant \alpha < s, \\[2pt]
    A_{\alpha - s}(X_1, \dots, X_{n-s})^{\,p^s} & \text{si } s \leqslant \alpha \leqslant n .
  \end{cases}
\]
\footnote{La page écrit ces formules sans les puissances de Frobenius :
$A_i(X_1, \dots, X_{n-1}, 0) = A_{i-1}(X_1, \dots, X_{n-1})$. Le twist est
nécessaire, et se lit directement sur la factorisation : en annulant $X_n$,
chaque valeur $\sum \lambda_i X_i$ est atteinte $p$ fois, donc
$\prod_{\lambda \in \mathbb{F}_p^n}(T - \lambda \cdot X) =
\bigl(\prod_{\mu \in \mathbb{F}_p^{n-1}}(T - \mu \cdot X)\bigr)^p$, et
élever au carré de Frobenius décale les exposants tout en élevant les
coefficients à la puissance $p$. Vérification pour $n = 2$ :
$A_1^{(2)}(X_1, 0) = -X_1^{p(p-1)}$, tandis que $A_0^{(1)}(X_1) =
-X_1^{p-1}$. La surjectivité n'en souffre pas : sur un corps algébriquement
clos, donc parfait, on extrait les racines $p^r$-ièmes qu'il faut.}

Avec ce lemme, on relève $(a_r, \dots, a_{n-1})$ par un $\xi'$ à $n - r$
variables — cas $r = 0$ en $n - r$ variables, après extraction des racines
$p^r$-ièmes — puis on complète par $r$ zéros.

\pagerange{11}{12}
\subsection*{La variante $SL$ (pages 11 et 12)}

Puisque $\Delta_n(u \cdot X) = \det(u)\, \Delta_n(X)$, le déterminant de
Moore est invariant sous $S\Gamma_n = SL(n, \mathbb{F}_p)$, et

\textbf{Corollaire.}
$\mathbb{F}_p[X_1, \dots, X_n]^{SL(n, \mathbb{F}_p)} =
\mathbb{F}_p[A_1, \dots, A_{n-1}, \Delta_n]$, avec
$\Gamma_n/S\Gamma_n \xrightarrow[\ \det\ ]{\ \sim\ } \mathbb{F}_p^*$ opérant
trivialement sur $A_1, \dots, A_{n-1}$ et par homothéties sur
$\Delta_n$.\footnote{$A_0$ n'a pas à figurer parmi les générateurs : il vaut
$(-1)^n \Delta_n^{p-1}$. La page barre d'ailleurs la ligne qui aurait dit que
les $\Delta_i = \Delta_n A_i$ sont eux aussi $S\Gamma_n$-invariants — ils le
sont, mais ils ne sont pas des générateurs indépendants.}

\emph{Démonstration.} Le carré
\[
  \begin{array}{ccc}
    \mathbb{A}^n/S\Gamma_n & \longrightarrow & \mathbb{A}^n \\
    \downarrow & & \downarrow \\
    \mathbb{A}^n/\Gamma_n & \xrightarrow{\ \sim\ } & \mathbb{A}^n
  \end{array}
\]
a pour flèche de droite celle définie par l'inclusion
$\mathbb{F}_p[A_1, \dots, A_{n-1}, \Delta_n^{p-1}] \subset
\mathbb{F}_p[A_1, \dots, A_{n-1}, \Delta_n]$. Les deux flèches verticales
sont des revêtements de groupe $\mathbb{F}_p^*$ identifiant le but au
quotient de la source ; la flèche du haut est compatible aux deux actions,
donc génériquement bijective et birationnelle, donc un isomorphisme puisque
$\mathbb{A}^n$ est normal.

\pagerange{13}{16}
\section*{La version relative : modules à Frobenius inversible (pages 13 à 16)}

Sur une base quelconque, il n'y a plus de « $n$ éléments indépendants » à
écrire en ligne. Ce qui les remplace est un faisceau, et toute la
construction des pages 7 à 12 se transpose — c'est le contenu de ces quatre
pages.

\pagerange{13}{13}
\subsection*{Structures spéciales linéaires et racines $(p-1)$-ièmes (page 13)}

\textbf{Corollaire.} Soit $M \subset \mathbb{G}_{a,S}$ fini localement libre
de rang $p^n$, d'où $a$ avec $M = \operatorname{Ker} f_a$ ; supposons $a_0$
inversible, c'est-à-dire $M$ étale. Alors l'ensemble des structures
$\mathbb{F}_p$-spéciales linéaires de $M$ — les trivialisations de
$\Lambda^n_{\mathbb{F}_p} M$ — est en correspondance biunivoque avec
l'ensemble des racines $(p-1)$-ièmes de $a_0$. Plus précisément
$(\Lambda^n_{\mathbb{F}_p} M)^*$ s'identifie au torseur, sous
$\mu_{p-1,S} \simeq (\mathbb{F}_p^*)_S$, de l'image réciproque de $a_0$ par
$\mathbb{G}_m \xrightarrow{\ x \mapsto x^{p-1}\ } \mathbb{G}_m$.

C'est la traduction exacte de $A_0 = (-1)^n \Delta_n^{p-1}$ : une racine
$(p-1)$-ième de $a_0$ est, au signe près, une valeur de $\Delta_n$,
c'est-à-dire un générateur de la puissance extérieure maximale, c'est-à-dire
une trivialisation.

\textbf{N.B.} Se donner sur $S$ un faisceau en $\mathbb{F}_p$-espaces
vectoriels de rang $n$, fini étale, équivaut à se donner un faisceau $E$
localement libre de rang $n$ et un isomorphisme
\[
  u : E^{(p)} \xrightarrow{\ \sim\ } E , \qquad E^{(p)} = F_S^* E ,
\]
la traduction étant $E = M \otimes_{\mathbb{F}_p} \mathcal{O}_S$ dans un
sens, et
$M = \operatorname{Ker}\bigl(V(E) \xrightarrow{\ \operatorname{id} -
u \circ F_E\ } V(E)\bigr)$ dans l'autre.\footnote{C'est l'équivalence entre
$\mathbb{F}_p$-systèmes locaux et \emph{modules à Frobenius inversible} —
aujourd'hui \emph{unit-root $F$-crystals} —, dont le fond est la théorie
d'Artin--Schreier : la suite exacte
$0 \to \mathbb{F}_p \to \mathbb{G}_a \xrightarrow{x^p - x} \mathbb{G}_a \to
0$ pour la topologie étale. La formulation qu'en donne la page est celle
qu'on trouve chez Katz ; rien sur le feuillet ne renvoie à une source, et on
n'affirme pas qu'il en avait une en vue.}

Se donner un homomorphisme $\varphi : M \to \mathbb{G}_a$ équivaut à se
donner un homomorphisme $\mathcal{O}$-linéaire $E \to \mathcal{O}_S$,
c'est-à-dire une section $\xi$ de $\check{E}$ ; et $\check{E}$ porte la
structure duale $\check{u} : \check{E}^{(p)} \to \check{E}$.

\pagerange{14}{15}
\subsection*{Le déterminant de Moore d'un faisceau (pages 14 et 15)}

Pour $\xi$ section de $\check{E}$, on pose
\[
  \xi^{(p)} \;=\; \check{u}\bigl(\xi^{(F)}\bigr) ,
\]
puis $\xi^{(p^i)}$ par itération ; ces sections correspondent aux composés de
$\varphi : M \to \mathbb{G}_a$ avec les itérées du Frobenius de
$\mathbb{G}_a$. On forme alors
\[
  \Delta_n^E(\xi) \;=\; \xi \wedge \xi^{(p)} \wedge \dots \wedge
  \xi^{(p^{n-1})} \;\in\; \Gamma\bigl(\Lambda^n \check{E}\bigr) ,
\]
soit une application polynomiale, homogène de degré
$\frac{p^n - 1}{p - 1}$,
\[
  \Delta_n^E : V(E) \longrightarrow V\bigl(\Lambda^n E\bigr) .
\]
L'image réciproque de la section nulle est un diviseur de Cartier relatif ;
son complémentaire $V(E)^0$ est un ouvert schématiquement dense. Sur lui,
$\xi, \xi^{(p)}, \dots, \xi^{(p^{n-1})}$ forment une base de $\check{E}$, et
l'on peut écrire
\[
  \xi^{(p^n)} + A^E_{n-1}(\xi)\,\xi^{(p^{n-1})} + \dots + A^E_0(\xi)\,\xi
  = 0 ,
\]
ce qui définit des fonctions
$A^E_i \in \Gamma(V(E)^0, \mathcal{O})$ uniquement déterminées. Comme dans le
cas absolu, elles s'étendent à $V(E)$ tout entier — les $A_i$ sont des
polynômes, pas seulement des fractions —, et l'on a encore
\[
  A^E_0 \;=\; (\pm 1)\,\bigl(\Delta_n^E\bigr)^{p-1} .
\]
Cette dernière égalité a un sens parce que
$(\Lambda^n \check{E})^{\otimes (p-1)} \simeq \mathcal{O}_S$ :
en effet $(\Lambda^n \check{E})^{\otimes p} \simeq
\Lambda^n(\check{E}^{(p)}) \xrightarrow{\ \Lambda^n \check{u}\ }
\Lambda^n \check{E}$.\footnote{C'est exactement l'argument entre crochets de
la page 15, et c'est ce qui rend la page 13 possible : la trivialisation de
$(\Lambda^n \check{E})^{\otimes(p-1)}$ est ce qui permet de parler d'un
torseur sous $\mu_{p-1}$ plutôt que d'une racine $(p-1)$-ième d'un scalaire.
Le signe, ici, on le laisse comme la page le laisse : il vaut $(-1)^n$ dans
la trivialisation du cas absolu, mais la trivialisation relative dépend de
$u$ et le feuillet ne la normalise pas.}

\pagerange{15}{16}
\subsection*{La matrice compagnon, et le tableau qui reste ouvert (pages 15 et 16)}

Si $\xi$ est une section de $V(E)^0$, les $\xi^{(p^i)}$ pour
$0 \leqslant i \leqslant n-1$ forment une base de $\check{E}$, qui s'identifie
ainsi à $\mathcal{O}_S^n$ ; en notant $e_i = \xi^{(p^i)}$ on a
\[
  e_0^{(p)} = e_1, \quad e_1^{(p)} = e_2, \quad \dots, \quad
  e_{n-2}^{(p)} = e_{n-1}, \qquad
  e_{n-1}^{(p)} = -\bigl(a_0 e_0 + a_1 e_1 + \dots + a_{n-1} e_{n-1}\bigr),
\]
c'est-à-dire que $\check{u}$ s'écrit, dans cette base, par la matrice
compagnon
\[
  \begin{pmatrix}
    0 & 0 & \cdots & 0 & -a_0 \\
    1 & 0 & \cdots & 0 & -a_1 \\
    0 & 1 & & \vdots & \vdots \\
    \vdots & & \ddots & 0 & \vdots \\
    0 & 0 & \cdots & 1 & -a_{n-1}
  \end{pmatrix} .
\]
\footnote{La dernière colonne du feuillet est en partie recouverte de
ratures ; la transcription ne garantit que $-a_0$, $a_1$ et $-a_{n-1}$. La
matrice compagnon est déterminée par les relations écrites juste au-dessus,
qui, elles, sont nettes, et c'est d'elles qu'on la tire — non de la lecture
des entrées.}

Le feuillet s'achève sur un tableau à deux colonnes, chacune de trois étages
reliés par des flèches montantes, les colonnes étant reliées par des
inclusions de gauche à droite.

La colonne de gauche est le cas étale, et ses trois étages sont équivalents :

\begin{itemize}
\item les $(a_0, \dots, a_{n-1})$ à $a_0$ inversible ;
\item les sous-groupes finis étales de $\mathbb{G}_{a,S}$ de rang $p^n$ ;
\item les couples formés d'un $E$ localement libre de rang $n$ avec
$\check{E}^{(p)} \simeq \check{E}$, et d'une section $\xi$ de $\check{E}$
telle que $\xi, \xi^{(p)}, \dots, \xi^{(p^{n-1})}$ soit une base.
\end{itemize}

La colonne de droite est le cas général, et c'est là que le feuillet
s'arrête :

\begin{itemize}
\item les $(a_0, \dots, a_{n-1})$ quelconques ;\footnote{Le feuillet écrit
$(a_0, \dots, a_n)$ en tête de cette colonne. Il en faut $n$, comme à gauche
et comme partout ailleurs dans le dossier ; on lit une inattention.}
\item les sous-groupes finis localement libres de rang $p^n$ de
$\mathbb{G}_{a,S}$ ;
\item les couples formés d'un $F$ localement libre de rang $n$ avec une
flèche $F^{(p)} \to F$ non nécessairement inversible, et d'une section $\xi$
de $F$ telle que $\xi, \xi^{(p)}, \dots, \xi^{(p^{n-1})}$ soit une base.
\end{itemize}

La flèche montante du troisième étage vers le deuxième, dans la colonne de
droite, porte un point d'interrogation, et une accolade court du bas vers le
haut de la colonne, annotée « trivial ? ». C'est le seul endroit du dossier
où l'auteur marque explicitement qu'il ne conclut pas.\footnote{On peut dire
ce qui est acquis et ce qui ne l'est pas. Du premier étage au troisième, il
n'y a rien à démontrer : la matrice compagnon ci-dessus fabrique
$(F, \phi, \xi)$ à partir de $a$, et la condition de base est vraie par
construction. Du premier au deuxième, c'est le théorème de la page 6. Ce qui
reste ouvert est ce que le point d'interrogation vise : une caractérisation
du deuxième étage par le troisième qui ne passe pas par le choix d'un $\xi$,
c'est-à-dire un énoncé de nature dieudonnéenne. Le feuillet ne dit pas
laquelle des deux choses il vise, et on ne tranche pas pour lui.}

\pagerange{18}{18}
\section*{Descente le long d'un pincement (page 18)}

La page 17 est blanche. La page 18 ouvre un sujet neuf et ne le referme pas.

On part d'un schéma $X'$ et de deux points $a, b$ de $X'$ ; $X$ est obtenu
en les identifiant — la marge dit « déduit de $X'$ en pinçant $a$ et
$b$ ».\footnote{C'est ce qu'on appelle aujourd'hui un \emph{pincement}, et
le formalisme complet de la descente le long d'un tel carré co-cartésien est
celui de Ferrand. Le feuillet n'y renvoie évidemment pas et n'en a pas
besoin : il manipule directement les données de recollement.} La question
est celle de la descente : si $H' \subset \mathbb{G}_{a,X'}$ est un
sous-groupe fini localement libre dont les fibres en $a$ et en $b$
coïncident, $H'_a = H'_b$, descend-il en un $H \subset \mathbb{G}_{a,X}$ ?

Le feuillet met en place l'obstruction plutôt que de répondre. Soit
$\varphi$ un homomorphisme de $\mathbb{G}_{a,X}$ s'annulant sur $H$, et $N'$
l'adhérence schématique dans $\mathbb{G}_{a,X'}$ du noyau de sa version
relevée, de sorte que
\[
  N' \subset H' \subset \mathbb{G}_{a,X'} , \qquad
  \mathbb{G}_{a,X'}/N' \longrightarrow \mathbb{G}_{a,X'}/H'
  \quad \text{étale} .
\]
Et il pose la situation qui l'intéresse :
\[
  H'(a) = H'(b) \qquad\text{mais}\qquad N'(a) \neq N'(b) .
\]
Autrement dit, la donnée de recollement peut valoir pour $H'$ sans valoir
pour le sous-groupe distingué $N'$ qu'on en a tiré : ce qui descend ne
descend pas avec sa filtration. Suivent deux suites exactes,
\[
  0 \to N' \to \mathbb{G}_{a,X'} \xrightarrow{\ f_{\alpha'}\ }
  \mathbb{G}_{a,X'} \to 0 , \qquad
  0 \to K' \to \mathbb{G}_{a,X'} \xrightarrow{\ f_{\beta'}\ }
  \mathbb{G}_{a,X'} \to 0 ,
\]
et des flèches $u_\lambda$, $u_{\lambda'}$ vers $\mathbb{G}_{a,X}$ annotées
« hom. Verschiebung », destinées à comparer les deux. L'argument qui devait
les comparer n'est pas écrit.\footnote{Ce n'est pas un défaut de lecture : la
transcription note que la fin de la page est une esquisse, les suites et les
flèches tracées, l'argument absent, et qu'un des points d'arrivée disparaît
sous une tache d'encre. On ne comble pas. Ce qui se tient — la mise en place
de l'obstruction et l'exemple que la page 22 en donnera — se tient sans lui.}

\pagerange{19}{22}
\section*{Le cas $n = 2$, calculé à la main (pages 19 à 22)}

Les quatre derniers feuillets sont du calcul, sans une phrase de rédaction
suivie. Ils ne répètent pas les pages précédentes : ils en font tourner la
machinerie sur le plus petit cas non trivial, et la page 22 revient à la
page 18 avec un exemple.

\pagerange{19}{20}
\subsection*{Les brouillons au crayon (pages 19 et 20)}

La page 19 réécrit le polynôme $X^{p^n} + A_1 X^{p^{n-1}} + \dots +
A_{n-1}X$, les $A_i$ y étant indexés en sens inverse de celui des pages 9 à
15.\footnote{On n'harmonise pas. L'indexation des pages 9 à 15 est celle où
$A_i$ est le coefficient de $X^{p^i}$, donc homogène de degré $p^n - p^i$ ;
celle de la page 19 compte à partir du haut. Les deux se traduisent l'une
dans l'autre par $i \mapsto n - i$, et la page 19 étant au crayon et sans
rédaction, il n'y a pas lieu de décider qu'elle corrige les précédentes.} Un
bloc écrit dans le sens de la hauteur occupe la marge gauche et reprend, au
crayon, le calcul d'espace tangent des pages 5 et 6 — $H_0$, $K_0 =
\mathbb{G}_0/H_0$, $\operatorname{Hom}(\alpha_{p^r}, \mathbb{G}_a)$, $F^r =
0$. Il est en grande partie illisible et n'ajoute rien qu'on puisse
établir.\footnote{Ce bloc est, avec les pages 1, 5-6 et 18, l'endroit du
dossier où la charge d'illisible est la plus lourde. On a renoncé à en tirer
autre chose que la constatation qu'il refait le calcul déjà fait ; il serait
facile, et faux, de le reconstruire depuis les pages 5 et 6.}

Suivent les relations de la matrice compagnon dans l'indexation renversée,
et un $M$ localement libre de type fini avec $M^{(p)} \xrightarrow{\ \sim\ }
M$ : c'est le module à Frobenius inversible des pages 13 à 15, rejoué au
crayon.

La page 20 essaie des coefficients pour $n = 2$, écrit le déterminant de
Moore sous forme transposée, redonne le degré $\frac{p^n-1}{p-1}$ et le
produit sur les formes linéaires — avec le même $\mathbb{P}^n$ pour
$\mathbb{P}^{n-1}$ qu'à la page 7. Deux lignes y méritent d'être relevées :
\[
  \frac{(X^p Y - Y^p X)^p}{X^p Y - Y^p X} = (X^p Y - Y^p X)^{p-1},
  \qquad \text{annoté } (X \wedge Y) ,
\]
qui est exactement $A_0 = \pm\Delta_2^{p-1}$ avec, dans la marge, la lecture
extérieure des pages 13 à 15 ; et
\[
  \left(\frac{X}{Y}\right)^{p^2} = \frac{X}{Y} ,
\]
c'est-à-dire que le rapport de deux éléments d'une base est fixé par le carré
du Frobenius.\footnote{La transcription signale que les exposants de la
première ligne de la page 20 sont les plus incertains du feuillet. Elle y
écrit $(f^{p-1} - g^{p-1})X^{p^2} - (f^{p^2} - g^{p^2})X^p + (f^{p^2}g^{p-1}
- g^{p^2}f^{p-1})X$, et cette ligne n'est pas une identité : la page 21
écrira la bonne, et on la donne là-bas. On ne corrige pas la page 20, dont
les exposants ne sont pas assurés.}

\pagerange{21}{21}
\subsection*{Élimination à partir de deux racines (page 21)}

Le feuillet part de $X^{p^2} + aX^p + bX$ et de deux racines $\alpha$,
$\beta$ :
\[
  \alpha^{p^2} + a\alpha^p + b\alpha = 0 , \qquad
  \beta^{p^2} + a\beta^p + b\beta = 0 .
\]
En multipliant la première par $\beta$, la seconde par $\alpha$ et en
retranchant, le terme en $b$ disparaît et
\[
  a\bigl(\alpha^p\beta - \beta^p\alpha\bigr) =
  -\bigl(\alpha^{p^2}\beta - \beta^{p^2}\alpha\bigr) ,
  \qquad\text{d'où}\qquad
  a = -\,\frac{\alpha^{p^2}\beta - \beta^{p^2}\alpha}
              {\alpha^p\beta - \beta^p\alpha} .
\]
De même, en multipliant par $\beta^p$ et $\alpha^p$,
\[
  b = \frac{\alpha^{p^2}\beta^p - \beta^{p^2}\alpha^p}
           {\alpha^p\beta - \beta^p\alpha} .
\]
\footnote{La page écrit $a$ sans le signe moins, et $b$ avec le bon signe. Le
dénominateur de $b$ porte d'ailleurs deux exposants $p$ que la page biffe
elle-même pour le ramener à celui de $a$ ; corrigé, il est exact. C'est
Cramer sur deux lignes, et c'est le cas $n = 2$ de la page 8.}

La colonne de droite recommence autrement, en utilisant $\alpha^p - \beta^p =
(\alpha - \beta)^p$ :
\[
  a(\alpha^p - \beta^p) = -(\alpha^{p^2} - \beta^{p^2}) ,
  \qquad\text{d'où}\qquad a = -(\alpha - \beta)^{p^2 - p} .
\]
\footnote{Cette première relation s'obtient en retranchant les deux
équations, mais la soustraction laisse un terme $b(\alpha - \beta)$ que la
page ne porte pas. La valeur $a = -(\alpha-\beta)^{p^2-p}$ qui en sort ne
coïncide donc pas avec celle de la colonne de gauche. On la signale sans la
réparer : tout ce qui suit dans cette colonne est exact \emph{à partir
d'elle}, et c'est la cohérence interne de la suite qui permet d'affirmer
qu'il s'agit d'un terme omis, non d'une lecture douteuse.}

De là, les trois lignes suivantes, qui sont exactes entre elles :
\[
  b\alpha = -\alpha^{p^2} + \alpha^p(\alpha - \beta)^{p^2 - p} ,
\]
\[
  (\alpha - \beta)^p\, b\alpha =
  \beta^p\alpha^{p^2} - \alpha^p\beta^{p^2}
  = \alpha^p\beta^p\bigl(\alpha^{p^2 - p} - \beta^{p^2 - p}\bigr) ,
\]
\[
  b\alpha = \Bigl(\frac{\alpha\beta}{\alpha - \beta}\Bigr)^{p}
  \bigl(\alpha^{p^2 - p} - \beta^{p^2 - p}\bigr) .
\]
\footnote{Les trois lignes sont écrites sur le feuillet avec $b$ au premier
membre là où c'est $b\alpha$ : la relation tirée de $\alpha^{p^2} + a\alpha^p
+ b\alpha = 0$ porte le facteur $\alpha$, et le même facteur manque aux deux
lignes suivantes, qui restent cohérentes entre elles. Rétabli, on obtient
$b = \alpha^{p-1}\beta^p(\alpha-\beta)^{-p}(\alpha^{p^2-p} -
\beta^{p^2-p})$. C'est la faute qui court sur trois lignes ; on la corrige
en une fois plutôt qu'à chaque ligne.}

Le feuillet dégage ensuite le polynôme obtenu en chassant les dénominateurs
et, après un détour par une version où les coefficients portent un facteur
$\lambda$ et sont nommés $r$, $s$, $t$, il finit par écrire :
\[
  \bigl(f^p g - g^p f\bigr) X^{p^2}
  - \bigl(f^{p^2} g - g^{p^2} f\bigr) X^{p}
  + \bigl(f^{p^2} g^{p} - g^{p^2} f^{p}\bigr) X .
\]
C'est, au signe près, le déterminant de Moore de $(f, g, X)$ :
\[
  \det\begin{pmatrix} f & g & X \\ f^p & g^p & X^p \\
  f^{p^2} & g^{p^2} & X^{p^2}\end{pmatrix} ,
\]
c'est-à-dire le $D(T)$ de la page 8 pour $n = 2$, dont les racines sont
exactement $\mathbb{F}_p f + \mathbb{F}_p g$. La dernière ligne du feuillet
rejoint donc la page 7, et c'est ce qui donne leur sens aux deux feuillets de
calcul : ils refont à la main, sur deux éléments, la construction que les
pages 7 à 16 mènent en général.\footnote{On l'a vérifié en développant le
déterminant selon la dernière colonne : les trois cofacteurs sont
$fg^p - gf^p$, $-(fg^{p^2} - gf^{p^2})$ et $f^pg^{p^2} - g^pf^{p^2}$, et la
ligne du feuillet en est l'opposée terme à terme. Les versions
intermédiaires de la même page, celles à coefficients $(\alpha^p - \beta^p)$
et celles à facteur $\lambda$, portent le défaut signalé plus haut ; la
dernière ligne ne le porte plus.}

\pagerange{22}{22}
\subsection*{Le feuillet écrit des deux bouts (page 22)}

La page 22 est écrite des deux bouts : le tiers supérieur à l'endroit, les
deux tiers inférieurs tête-bêche, séparés par un long trait au crayon, le
bloc renversé se lisant depuis le pied du feuillet. Ce n'est pas un défaut de
numérisation mais la mise en page de l'auteur, et elle fixe l'ordre de
lecture ; sous son trait, le bloc renversé est lui-même en deux colonnes.

\emph{Le bloc à l'endroit} construit l'exemple que la page 18 appelait. On y
prend deux fonctions $f$, $g$ sur $X'$ avec
\[
  f(a) = 0, \quad g(a) = 1, \qquad f(b) = 1, \quad g(b) = 0 ,
\]
donc non constantes et linéairement indépendantes sur $\mathbb{F}_p$. À
chacune est associé $H_f, H_g \subset \mathbb{G}_a$, et on pose $H = H_f
\times H_g$. Le point est alors :
\[
  H_a \simeq \alpha_p \times \mathbb{Z}/p , \qquad
  H_b \simeq \mathbb{Z}/p \times \alpha_p .
\]
Les deux fibres sont abstraitement isomorphes, mais les deux facteurs se
croisent entre $a$ et $b$ : l'isomorphisme existe, la décomposition ne se
recolle pas. C'est exactement ce que la page 18 avait besoin de produire — un
$H'$ dont la donnée de recollement vaut, et dont les sous-groupes distingués
ne la respectent pas.

Le feuillet se donne ensuite $\lambda$ non constante avec
$\lambda(a) = \lambda(b) = 0$ — c'est-à-dire une fonction qui, elle, descend
au pincement.

\emph{Le bloc renversé} travaille l'opérateur qui fabrique $H_f$. On y lit,
en regard l'une de l'autre, les deux écritures
\[
  \Bigl(\frac{\lambda}{f}\Bigr)^{p} - \frac{\lambda}{f}
  \qquad\text{et}\qquad
  \lambda^p - f^{p-1}\lambda ,
\]
qui diffèrent du facteur $f^p$. Le second membre est le polynôme additif
$\wp_f(\lambda) = \lambda^p - f^{p-1}\lambda$, dont le lieu de zéros est
exactement la droite $\mathbb{F}_p f$ : c'est un Artin--Schreier tordu par
$f$, et c'est le cas $n = 1$ de la construction de la page 7, puisque
$D(T)/\Delta_1 = T^p - f^{p-1}T$.

Le feuillet compose alors deux opérateurs de cette espèce :
\[
  r\bigl(\wp_f \lambda\bigr)^{p} - r^{p}\bigl(\wp_f \lambda\bigr) ,
  \qquad r = g^p - f^{p-1}g = \wp_f(g) ,
\]
c'est-à-dire $r \cdot \wp_r(\wp_f(\lambda))$. Le lieu de zéros de ce composé
est $\mathbb{F}_p f + \mathbb{F}_p g$ : c'est, à un scalaire près, le
polynôme de degré $p^2$ que la dernière ligne de la page 21 avait écrit
comme un déterminant.\footnote{La transcription signale que $r$ est d'abord
écrit $(g/f)^p - f^{p-1}g$, le quotient $g/f$ étant biffé d'une croix ;
corrigé, $r = g^p - f^{p-1}g$, qui est bien $\wp_f(g)$. On note au passage
que $f\,\wp_f(g) = fg^p - f^pg = \Delta_2(f,g)$ : le coefficient dominant du
composé est le déterminant de Moore des deux fonctions, divisé par $f$. Les
deux présentations — composée d'Artin--Schreier tordus, et quotient de
déterminants — sont la même chose vue des deux bouts, ce qui est aussi le
cas des pages 7 à 9 et des pages 13 à 15.} La ligne suivante,
\[
  \left(\frac{\bigl(\frac{\lambda}{f}\bigr)^p - \frac{\lambda}{f}}
             {r}\right)^{\nu} = (\quad) ,
\]
amorce visiblement l'itération suivante — on renormalise par $r$ avant de
réappliquer l'opérateur — mais l'exposant est un seul signe, lu $\nu$, et le
second membre est une parenthèse laissée vide. On n'en tire rien.

La suite du bloc renversé reprend la situation à deux points sous une forme
un peu plus riche : $H = H_1 \times H_2 \times K$ avec
$H_{1a} \simeq H_{2b} \simeq \alpha_p$ — le croisement des facteurs, de
nouveau —, puis $f(a) = 0$, $g(b) = 0$, $f(b) = \alpha$, $g(a) = \beta$ avec
$\alpha, \beta$ indépendants sur $\mathbb{F}_p$, et la fibre
\[
  \alpha_p + \mathbb{F}_p\alpha + \mathbb{F}_p\beta \subset \mathbb{G}_a ,
\]
de rang $p^3$. La colonne de gauche écrit alors la donnée de recollement
elle-même :
\[
  i_a(H'_a) = i_b(H'_b) = \alpha_p + \mathbb{F}_p\alpha +
  \mathbb{F}_p\beta ,
\]
c'est-à-dire que les deux fibres coïncident comme sous-groupes de
$\mathbb{G}_a$, ce qui est précisément la condition de descente de la
page 18 ; et le $H \subset \mathbb{G}_{a,X}$ correspondant a la même fibre.
La colonne de droite ajoute une troisième fonction $h$, avec
$h(a), h(b) \in \mathbb{F}_p\alpha + \mathbb{F}_p\beta$,
$h(a) \notin \mathbb{F}_p\beta$ et $h(b) \notin \mathbb{F}_p\alpha$ — par
exemple $h(a) = \alpha$, $h(b) = \beta$ —, ce qui est la façon de rendre le
croisement inévitable.

Le feuillet finit sur un diagramme où $H \hookrightarrow \mathbb{G}_a
\xrightarrow{\ f\ } \mathbb{G}_a$ est complété par une flèche $u$ et une
flèche verticale qui porte un point d'interrogation, puis sur
$\mathbb{G}_{a,X'}/H_1 \times H_2$ et, au crayon fin, sur
$\lambda \mapsto (\lambda/f)^p - (\lambda/f)$. C'est le second et dernier
point d'interrogation du dossier ; le premier était sur la colonne droite du
tableau de la page 16.\footnote{Plusieurs annotations de ce bloc ne se lisent
pas : deux mots au crayon sous une accolade, deux lignes biffées dont ne
surnage que $H_2$, trois lignes entre accolades à droite de $f$, $g$. Elles
portaient sans doute l'énoncé que la page voulait atteindre. On ne le
devine pas.}

\section*{Ce que le dossier annonce et n'établit pas}

Il vaut la peine de les rassembler, parce qu'aucun n'est dû à l'état du
papier :

\begin{itemize}
\item \textbf{La réduction de $X_n$.} Le théorème de la page 6 est énoncé
dans sa généralité ; ce que les pages 2 à 5 démontrent est l'isomorphisme sur
$X_{n,\mathrm{réd}}$, et la question de la réduction, explicitement posée à
la page 5, n'est traitée qu'au cas $n = 1$ et au seul point $0$.
\item \textbf{La constante $c$} de la factorisation du déterminant de Moore :
« reste à déterminer $c$ », page 7, et elle n'est pas déterminée. Elle vaut
$1$ pour la normalisation que la page adopte.
\item \textbf{L'ouverture de $X_n^{\text{ét}}$}, portée « à justifier » en
marge de la page 4.
\item \textbf{La colonne de droite du tableau de la page 16}, avec son point
d'interrogation et son « trivial ? ».
\item \textbf{La descente de la page 18}, dont la mise en place est écrite et
l'argument non.
\item \textbf{L'itération de la page 22}, qui s'arrête sur une parenthèse
vide.
\end{itemize}

Et une dernière chose, qui n'est pas une lacune mais un fait : le lemme de la
page 1 n'a de rapport avec rien d'autre dans le dossier. Ni sa question — la
densité schématique d'une section formelle le long d'un fermé — ni son
vocabulaire ne reparaissent aux vingt et une pages suivantes. C'est un
feuillet qui s'est trouvé là.

\end{document}
